
\lussec 5. Lattice regularization

The starting point in this section is a formulation of lattice
QCD, where the quark fields are put on the lattice as proposed
by Wilson [\cite{Wilson}]. 
No particular assumptions
are made on the lattice Dirac operator or the gauge action,
but both should be local and respect all symmetries preserved
by the standard Wilson theory. For notational convenience,
the lattice spacing $a$ is set to unity.
The lattice is taken to be infinitely extended in all directions 
unless specified otherwise.


\lussubsec 5.1 Flow equations

On the lattice, the time-dependent gauge field $V(t,x,\mu)$ is
defined by 
\equation{
  \left.V\right|_{t=0}=U,
  \enum
  \nexteq{2.5ex}
  \partial_tVV^{-1}=-g_0^2\partial\Sw,
  \enum
}
where $U(x,\mu)$ is the fundamental lattice gauge field 
and $(\partial\Sw)(V,x,\mu)$ the gradient of the
Wilson plaquette action (see appendix A and
ref.~[\cite{WilsonFlow}] for further explanations).

In the case of the quark fields, the boundary conditions 
and the form of the flow equations are the same as 
in the continuum theory, but the operator $\Delta$
in eq.~(2.4) gets replaced by the lattice Laplacian
\equation{
  \Delta=\nabstar{\mu}\nab{\mu},
  \enum
}
$\nab{\mu}$ and $\nabstar{\mu}$
being the gauge-covariant forward and backward
difference operators in presence of the time-dependent gauge field.

As already noted in ref.~[\cite{WilsonFlow}],
the global existence, uniqueness and smoothness of the solution
of the flow equation (5.2) is rigorously guaranteed.
Since $\Delta$ is a bounded operator, and since
it depends smoothly on the flow time, the 
same can be shown to be true in the case of quark flow equations
[\cite{ArnoldI}].
Local fields at non-zero flow time, such as those considered
in the previous sections, are thus completely well defined
on the lattice.


\lussubsec 5.2 Lattice theory in 4+1 dimensions

Similarly to the continuum theory,
the lattice QCD correlation functions of the 
time-dependent local fields can be obtained
from a lattice field theory in 4+1 dimensions.
To be able to fully control this construction,
the latter is first set up for a discretized
version of the flow equations, where the flow time is restricted
to a finite set
\equation{
  0,\eps,2\eps,3\eps,\ldots,T
  \enum
}
of values
separated by a time step $\eps>0$. 

The integration of the discretized evolution equation 
\equation{
  V(t+\eps,x,\mu)V(t,x,\mu)^{-1}=\rme^{-\eps g_0^2(\partial\Sw)(V,x,\mu)}
  \enum
}
for the gauge field amounts to an approximate integration of the 
continuous equation (5.2) using the Euler method
(or, equivalently, through the application
of a sequence of ``stout smearing'' steps [\cite{Stout}]).
In the case
of the quark fields, the discretized flow equation has the 
same form as the continuous one, but the time derivative 
of the quark field is replaced by the forward difference
\equation{
  \partial_t\chi(t,x)={1\over\eps}\{\chi(t+\eps,x)-\chi(t,x)\}.
  \enum
}
Note that the discrete equations are such that they can be
solved in steps,
starting from the fundamental fields at time $t=0$
and recursively proceeding from time $t$ to $t+\eps$.
In particular, the solutions are uniquely determined by the
initial values of the fields.

In the lattice theory in 4+1 dimensions, all fields integrated
over in the functional integral are defined at the times (5.4)
only. 
Since gauge fixing is not needed in the present context,
there are no ghost fields and no gauge terms in the action,
but the boundary conditions (2.6) and (5.1) are imposed 
as before.
The components of the Lagrange-multiplier fields at time $T$ 
decouple from the other fields and can therefore be set
to zero (or be omitted from the beginning). 

A possible choice
of the bulk actions is then
\equation{
  \SGfl=-2\sum_{t=0}^{T-\eps}\sum_{x,\mu}
  \tr\bigl\{L(t,x,\mu)
  \noenum
  \nexteq{-0.5ex}
  {\phantom{\SGfl=-2\sum_{t=0}^{T-\eps}\sum_{x,\mu}}}
  \times
  \Psu\bigl(
  V(t+\eps,x,\mu)V(t,x,\mu)^{-1}-\rme^{-\eps g_0^2(\partial\Sw)(V,x,\mu)}
  \bigr)\bigr\},
  \enum
  \nexteq{2.5ex}
  \SFfl=\eps\sum_{t=0}^{T-\eps}\sum_x\bigl\{
  \lambdabar(t,x)\left(\partial_t-\Delta\right)\chi(t,x)
  +\chibar(t,x)\left(\lvec{\partial}_t-\lvec{\Delta}\right)\lambda(t,x)
  \bigr\},
  \enum
}
where the function
\equation{
  M\to\Psu(M)={1\over2}(M-M^{\dagger})
  -{1\over2N}\kern1pt\tr(M-M^{\dagger})
  \enum
}
maps any complex $N\times N$ matrix to an element of the Lie algebra
of $\SUn$. Apart from these actions and the lattice QCD action,
the measure term
\equation{
  \Sms=\sum_{t=0}^{T-\eps}\sum_{x,\mu}
  {\cal L}_{\rm ms}(V(t+\eps,x,\mu)V(t,x,\mu)^{-1}),
  \enum
  \nexteq{3.0ex}
  {\cal L}_{\rm ms}(W)=
  \cases{-\ln\det(J(W)) & if $W\in\Dsu$,\cr
         \noalign{\vskip1ex}
         \infty         & otherwise,\cr}
  \enum
}
must be included in the total action of the lattice theory (see appendix B
for the definition of the Jacobian matrix $J(W)^{ab}$ of the mapping (5.9) 
and the neighbourhood $\Dsu\subset\SUn$ of unity). 

Having specified the fields and their action,
the functional integral of the lattice theory in 4+1 dimensions
is defined in the standard manner. 
Note that the measure term (5.10),(5.11) 
effectively restricts the link variables 
in the functional integral to the domain where
\equation{
  V(t+\eps,x,\mu)V(t,x,\mu)^{-1}\in\Dsu
  \quad\hbox{for all $t,x,\mu$}.
  \enum
}
Since the action $\SGfl$ is purely imaginary, 
the integral over the Lagrange-multiplier field $L(t,x,\mu)$ 
is not absolutely convergent, but it is understood
that the integration ambiguity this may entail is removed 
through the addition of an
infinitesimal quadratic term in the Lagrange-multiplier field
to the action.


\lussubsec 5.3 Recovering the theory in four dimensions

The correlation functions defined in this way of local fields 
composed from the fields $V$, $\chi$ and $\chibar$
can be evaluated, to some extent,
by first integrating over the Lagrange-multiplier fields.
Since the total action is linear in the latter, 
the integration yields a product
\equation{
  \prod_{t=0}^{T-\eps}\prod_{x,\mu}
  \delta\bigl\{\Psu\bigl(
  V(t+\eps,x,\mu)V(t,x,\mu)^{-1}\bigr)-
  \Psu\bigl(\rme^{-\eps g_0^2(\partial\Sw)(V,x,\mu)}\bigr)\bigr\}
  \noenum
  \nexteq{1.0ex}
  \times\prod_{t=0}^{T-\eps}\prod_{x}
  \delta\bigl\{\partial_t\chi(t,x)-\Delta\chi(t,x)\bigr\}
  \delta\bigl\{\chibar(t,x)\lvec{\partial}_t-
  \chibar(t,x)\lvec{\Delta}\bigr\}
  \enum
}
of Dirac $\delta$-functions.
The arguments of these 
$\delta$-functions vanish if $V$, $\chi$ and $\chibar$ satisfy the discrete
flow equations, and for sufficiently small step sizes $\eps$
there is in fact no other way the arguments 
can be made to vanish%
\kern1pt\sfn{$\dagger$}{%
Uniqueness is guaranteed if 
$\exp\{\eps g_0^2(\partial\Sw)(V,x,\mu)\}\in\Dsu$
on all links of the lattice and at all flow times.
Some rigorous estimates show that
this condition is satisfied
if $\eps\leq(24\sqrt{2}N)^{-1}$.}.
In the case of the gauge field,
the important point to note here is
that the mapping $\Psu$ in the $\delta$-functions 
operates on $\SUn$ matrices in a small neighbourhood of unity.
The properties of the mapping quoted in appendix B then
imply that the constraints imposed by the $\delta$-functions 
are equivalent to imposing the flow equation (5.5) on the gauge field.

The integral over the field variables at all flow times $t>0$ 
can now be performed using the $\delta$-functions. Starting
with the fields at time $t=T$ and proceeding from there 
to the smaller times, the $\delta$-functions allow 
all fields to be eliminated recursively until the only
integration variables left are the fundamental fields $U$, 
$\psi$ and $\psibar$. In this process, the integrals over the 
$\delta$-functions for the gauge field variables give rise
to a product of Jacobians of the mapping $\Psu$, but 
these factors are canceled by the measure term in the action
(see appendix B).

The calculation thus shows that the correlation functions
in 4+1 dimensions of local fields that do not involve the 
Lagrange-multiplier fields
are exactly equal to the correlation functions
of the same fields
defined in 4 dimensions through the discrete flow 
equations and the QCD functional integral.


\lussubsec 5.4 Continuous time limit

At this point, the time-dependent fields in the correlation functions
are still the ones obtained by the Euler integration of the flow equations.
Since the associated
integration errors are of order $\eps$,
the continuous-time correlation functions are recovered 
when $\eps\to0$ (and, trivially, $T\to\infty$).
Moreover, the measure term
\equation{
  \Sms=\sum_{t=0}^{T-\eps}\sum_{x,\mu}
  \ln\det J\bigl(\rme^{-\eps g_0^2(\partial\Sw)(V,x,\mu)}\bigr)
  =\rmO(\eps)
  \enum
}
vanishes in this limit. 

In the following, the continuous-time formulation of the lattice
theory will normally be considered, but 
the functional integral in 4+1 dimensions is always assumed to be
defined through the limit $\eps\to0$ of
the discrete-time functional integral.
